% =============================================================================
% 8. CONCLUSION
% =============================================================================
\section{Conclusion}
\label{sec:conclusion}

This study set out to investigate how the Kirkpatrick four-level model can
be operationalised within the training process of an \AS{}-certified
aircraft-engine \MRO{} facility so that it evaluates not only safety and
regulatory compliance but also resource efficiency and waste reduction
(RQ1), what measurable changes emerge from such an implementation over a
twelve-month horizon (RQ2), and what managerial and policy interventions
can be derived (RQ3).

\subsection{Answers to the research questions}
\label{subsec:rq-answers}

\paragraph{RQ1.} Kirkpatrick can be operationalised in an \AS{} \MRO{}
facility by extending the existing training process with three new
activities (evaluation across the four levels, verification against
acceptance criteria, data-driven optimisation), without disrupting the
underlying \LMS{}--\ERP{} architecture. Each level maps onto a specific
cluster of \AS{} clauses and \SMS{} pillars
(Fig.~\ref{fig:conceptual}).

\paragraph{RQ2.} Over twelve months, the intervention was associated
with a 28\% reduction in training failures, a 35\% reduction in
regulatory non-conformances, an approximately 6\% reduction in turnaround
time and an across-the-board improvement in customer satisfaction. The
associated annual quality-cost savings are estimated at USD~0.75--1.5~M
across pessimistic-to-optimistic scenarios (USD~1.19~M base case), with a
payback period below twelve months in every scenario. The case-study
design does not isolate the contribution of training from concurrent
operational changes; the findings are therefore reported as associations
rather than as causal effects (Section~\ref{subsubsec:validity}).

\paragraph{RQ3.} Three classes of intervention are recommended:
(i) at the facility level, embedding the four-level evaluation in monthly
quality reviews and gating \ERP{} access on Level~2 success;
(ii) at the industry level, joint IATA/SAE guidance on
training-evaluation methodology; (iii) at the regulator level, mandating
Level~3 and Level~4 reporting as part of \SMS{} safety performance
indicators.

\subsection{Contribution to the literature}
\label{subsec:contribution-lit}

The study contributes to three converging literatures. To the literature
on training evaluation, it provides --- to the best of the authors'
knowledge --- the first documented twelve-month operationalisation of
Kirkpatrick that jointly maps the four levels to \AS{} clauses, \SMS{}
pillars and \ESG{} indicators in an aerospace \MRO{} facility, with
monetised cost-of-quality scenarios. To the literature on aviation
safety management, it adds evidence that systematic training evaluation
is associated with reductions in regulatory non-conformances and quality
costs at a magnitude comparable to recent benchmarks
\citep{Zhang2025,Sousa2025}. To the literature on cleaner production and
circular economy, it adds a high-leverage example of how human-capital
interventions translate, under conservative assumptions, into avoided
material, energy and labour waste.

\subsection{Implications for practice}
\label{subsec:impl-practice}

For \MRO{} managers, the study offers a five-step roll-out checklist
(Section~\ref{subsec:mgr-integration}) and an audit-readiness map against
\AS{} clause~7.2 (Section~\ref{subsec:mgr-audit}). For airline quality
departments, it offers a defensible link between training metrics and
\ESG{} disclosures. For regulators, it offers a concrete grammar for
extending \SMS{} acceptance criteria to include training-effectiveness
reporting.

% =============================================================================
% Limitations and future research -- now embedded as subsection of Conclusion
% =============================================================================
\subsection{Limitations and future research}
\label{subsec:limitations}

The study has acknowledged limitations that bound the generalisation of its
findings. Table~\ref{tab:limitations} consolidates each limitation, the
mitigation adopted in the present study and the corresponding future-research
direction.

% =============================================================================
% TABLE 2 -- Limitations and future research  (Reviewer comment #6)
% =============================================================================
\begin{table*}[t]
  \centering
  \caption{Limitations of the present study, mitigation strategies and
           future research directions.}
  \label{tab:limitations}
  \small
  \begin{tabularx}{\textwidth}{@{}p{3.5cm} X X@{}}
    \toprule
    \textbf{Limitation} & \textbf{Mitigation in the present study} &
    \textbf{Future research direction} \\
    \midrule
    Single-case study (one \MRO{} facility) &
    In-depth triangulation across four data sources (questionnaires,
    observations, operational records, customer surveys). &
    Multi-site replication across geographic regions and regulatory regimes
    (\FAA{}, \EASA{}, \ANAC{}). \\
    \addlinespace
    Inability to fully isolate training effects from other variables &
    Workforce size and engine programmes held constant where feasible;
    \ROI{} sensitivity bounds reported. &
    Longitudinal quasi-experimental or causal-AI attribution
    \citep{Fernandez2024,Patel2023} to disentangle contemporaneous
    confounders. \\
    \addlinespace
    Self-reported behaviour data &
    Behaviour observations complemented by objective quality and
    customer-satisfaction metrics. &
    IoT-based behavioural tracking and digital-twin telemetry
    \citep{Garcia2024} for continuous Level~3 measurement. \\
    \addlinespace
    Generalisability to smaller \MRO{} facilities &
    Five-step roll-out checklist designed for non-disruptive integration
    into existing \LMS{}/\ERP{} infrastructure. &
    Adapted framework for low-resource settings (simplified Level~3
    observation protocol, sampled Level~4 metrics). \\
    \addlinespace
    Sustainability outcomes reported as order-of-magnitude rather than
    fully quantified &
    Cost-of-quality decomposition provides bounded estimates of avoided
    material, labour and energy. &
    Joint analysis with life-cycle assessment (LCA) and Scope~3 emissions
    accounting to translate \CoQ{} into specific environmental indicators. \\
    \addlinespace
    Aggregated data only (no raw individual records released) &
    Transparency about anonymisation and aggregation in figure captions
    and methodology; aggregated statistics reported by the facility's
    quality department. &
    Multi-organisation consortium for shared anonymised data sets that
    support cross-validation of the framework. \\
    \addlinespace
    Insider access -- one author is affiliated with the host organisation,
    which yields privileged data but may introduce framing or
    organisational bias &
    Inclusion of objective external metrics (regulatory audit outcomes,
    independent customer satisfaction surveys) that are verifiable
    outside the host organisation; explicit declaration in the competing
    interest statement. &
    External replication by independent academic groups in unaffiliated
    \MRO{} facilities to triangulate the present findings. \\
    \bottomrule
  \end{tabularx}
\end{table*}


In addition to the directions listed in the table, three broader research
avenues are noteworthy. First, comparing the framework across regulatory
regimes (\FAA{} versus \EASA{} versus \ANAC{}) would test whether the
mapping of Kirkpatrick levels to regulatory clauses requires
jurisdiction-specific adaptation. Second, building on the host facility's
parallel digital-twin programme for engine-specific procedural rehearsal
\citep{Garcia2024}, future work will examine whether digital-twin-enabled
Level~3 measurement provides continuous and unobtrusive behavioural data,
which would mitigate the Hawthorne-effect concern identified in
Section~\ref{subsubsec:validity}. Third, formal causal-attribution methods
\citep{Fernandez2024} could be applied to longitudinal data sets to
disentangle the contribution of training from contemporaneous changes in
workforce composition, tooling and maintenance volume.

